\documentclass[11pt,a4paper]{article}

% ── Codificación y fuentes ────────────────────────────────────────────────────
\usepackage[utf8]{inputenc}
\usepackage[T1]{fontenc}
\usepackage{lmodern}
\usepackage[spanish]{babel}

% ── Geometría y espaciado ─────────────────────────────────────────────────────
\usepackage[top=2.5cm, bottom=2.5cm, left=2.8cm, right=2.8cm]{geometry}
\usepackage{setspace}
\setstretch{1.15}
\usepackage{parskip}

% ── Tablas ────────────────────────────────────────────────────────────────────
\usepackage{booktabs}
\usepackage{tabularx}
\usepackage{array}
\usepackage{longtable}
\usepackage{enumitem}

% ── Colores y cajas ───────────────────────────────────────────────────────────
\usepackage[dvipsnames]{xcolor}
\usepackage{tcolorbox}
\tcbuselibrary{skins, breakable}

% ── Cabeceras y pies de página ────────────────────────────────────────────────
\usepackage{fancyhdr}
\usepackage{lastpage}
\pagestyle{fancy}
\fancyhf{}
\fancyhead[L]{\small\textbf{Informe de Evaluación} --- Física para Bioanalistas}
\fancyhead[R]{\small Daniela Vallenilla}
\fancyfoot[C]{\small Página \thepage\ de \pageref{LastPage}}
\renewcommand{\headrulewidth}{0.4pt}

% ── Hiperenlaces ──────────────────────────────────────────────────────────────
\usepackage[colorlinks=true, linkcolor=NavyBlue, urlcolor=NavyBlue]{hyperref}

% ── Paleta de colores personalizados ─────────────────────────────────────────
\definecolor{desempeno_excelente}{HTML}{1A6B2F}
\definecolor{desempeno_bueno}{HTML}{2E5A9C}
\definecolor{desempeno_suficiente}{HTML}{8B6914}
\definecolor{desempeno_deficiente}{HTML}{C0392B}
\definecolor{desempeno_insuficiente}{HTML}{7B241C}
\definecolor{grisFondo}{HTML}{F5F5F5}
\definecolor{azulTitulo}{HTML}{1B2A6B}
\definecolor{verdeFortaleza}{HTML}{1A6B2F}
\definecolor{naranjaMejora}{HTML}{A04000}

% ── Estilos de cajas ─────────────────────────────────────────────────────────
\tcbset{
  cajaTitulo/.style={
    enhanced, breakable,
    colback=azulTitulo!8, colframe=azulTitulo,
    fonttitle=\bfseries\large, coltitle=white,
    attach boxed title to top left={yshift=-2mm, xshift=4mm},
    boxed title style={colback=azulTitulo},
    top=6pt, bottom=6pt, left=8pt, right=8pt,
  },
  cajaFortaleza/.style={
    enhanced, breakable,
    colback=verdeFortaleza!6, colframe=verdeFortaleza!60,
    fonttitle=\bfseries, coltitle=verdeFortaleza,
    top=4pt, bottom=4pt, left=8pt, right=8pt,
  },
  cajaMejora/.style={
    enhanced, breakable,
    colback=naranjaMejora!6, colframe=naranjaMejora!60,
    fonttitle=\bfseries, coltitle=naranjaMejora,
    top=4pt, bottom=4pt, left=8pt, right=8pt,
  },
  cajaRecomendacion/.style={
    enhanced, breakable,
    colback=grisFondo, colframe=gray!50,
    fonttitle=\bfseries, coltitle=black,
    top=4pt, bottom=4pt, left=8pt, right=8pt,
  },
}

% ─────────────────────────────────────────────────────────────────────────────
\begin{document}

% ══ PORTADA ══════════════════════════════════════════════════════════════════
\begin{center}
  {\Large\bfseries\color{azulTitulo} Universidad de Oriente/Núcleo Sucre --- Física para Bioanalistas}\\[4pt]
  {\large Informe de Evaluación Preliminar}\\[12pt]
  \rule{\linewidth}{1.2pt}\\[6pt]
  {\LARGE\bfseries Daniela Vallenilla}\\[4pt]
  \rule{\linewidth}{0.6pt}
\end{center}

% ══ FICHA TÉCNICA ════════════════════════════════════════════════════════════
\vspace{4pt}
\begin{tcolorbox}[cajaTitulo, title=Datos del informe]
\begin{tabular}{@{}llll@{}}
  \textbf{Estudiante:}  & Daniela Vallenilla  &
  \textbf{Fecha:}       & 2026-08-08 \\[4pt]
  \textbf{Archivo PDF:} & \texttt{Daniela\_Vallenilla.pdf}  &
  \textbf{Páginas:}     & 6 \\
\end{tabular}
\end{tcolorbox}

% ══ RESULTADO GLOBAL ═════════════════════════════════════════════════════════
\vspace{6pt}
\begin{center}
  \begin{tcolorbox}[
    enhanced,
    colback=white, colframe=desempeno_bueno,
    width=0.62\linewidth,
    boxrule=2pt,
    halign=center, valign=center,
    top=8pt, bottom=8pt,
  ]
    \centering
    {\large\bfseries Puntaje sugerido}\\[6pt]
    {\Huge\bfseries\color{desempeno_bueno} 7.5/10}\\[6pt]
    {\large Nivel de desempeño: \textbf{\color{desempeno_bueno} Bueno}}
  \end{tcolorbox}
\end{center}

% ══ RESUMEN GENERAL ══════════════════════════════════════════════════════════
\section*{Resumen general}
El estudiante demuestra un buen entendimiento de los principios físicos fundamentales aplicados a fenómenos biológicos y clínicos. Las respuestas cualitativas son generalmente claras y precisas, especialmente en las preguntas 1 y 2. Sin embargo, se observan deficiencias en la atención al detalle en las conversiones de unidades, la omisión de partes de preguntas y errores en cálculos numéricos o la falta de procedimiento en algunas secciones.

% ══ FORTALEZAS Y ÁREAS DE MEJORA ════════════════════════════════════════════
\vspace{4pt}
\begin{minipage}[t]{0.48\linewidth}
  \begin{tcolorbox}[cajaFortaleza, title={Fortalezas}]
\begin{itemize}[leftmargin=1.5em, itemsep=2pt]
  \item Excelente comprensión conceptual de la termorregulación y la eficiencia de la sudoración (Pregunta 1).
  \item Dominio claro de la mecánica respiratoria y la Ley de Laplace, incluyendo el papel del surfactante (Pregunta 2).
  \item Capacidad para aplicar la Ley de Fick y analizar la eficiencia de difusión (Pregunta 3a).
  \item Correcta explicación del fenómeno de hemólisis en soluciones hipotónicas (Pregunta 4a).
  \item Habilidad para realizar cálculos numéricos cuando las fórmulas y datos son aplicados correctamente (Pregunta 1c, 2c, 4c).
\end{itemize}
  \end{tcolorbox}
\end{minipage}
\hfill
\begin{minipage}[t]{0.48\linewidth}
  \begin{tcolorbox}[cajaMejora, title={Areas de mejora}]
\begin{itemize}[leftmargin=1.5em, itemsep=2pt]
  \item Precisión en las conversiones de unidades en los cálculos.
  \item Asegurarse de responder todas las partes de una pregunta, incluyendo las explicaciones cualitativas.
  \item Mostrar el procedimiento completo de los cálculos, incluyendo fórmulas y sustitución de valores.
  \item Mejorar la exactitud en los cálculos numéricos.
\end{itemize}
  \end{tcolorbox}
\end{minipage}

% ══ OBSERVACIONES POR PREGUNTA ═══════════════════════════════════════════════
\section*{Observaciones por pregunta}

\begin{longtable}{>{\bfseries}p{2.8cm} p{1.6cm} p{7.0cm} p{2.8cm}}
  \toprule
  \textbf{Pregunta} & \textbf{Puntaje} & \textbf{Evaluación} & \textbf{Errores detectados} \\
  \midrule
  \endhead
  \bottomrule
  \endfoot
    \textbf{Pregunta 1: Termorregulación y Eficiencia de la Sudoración} & 2.0/2.0 & El estudiante responde de manera excelente a todas las partes de la pregunta. La justificación de la mayor sudoración en el paciente B es clara y precisa, explicando correctamente el rol del calor latente de vaporización. La explicación del efecto de la humedad relativa del 100\% es completa y bien fundamentada. Los cálculos de calor disipado en Joules y Kilocalorías son correctos y muestran las unidades adecuadas. & \textit{---} \\
    \midrule
    \textbf{Pregunta 2: Mecánica Respiratoria y Ley de Laplace} & 2.0/2.0 & La respuesta es sobresaliente. La explicación de la Ley de Laplace y el flujo anómalo de aire entre alvéolos de diferente tamaño es muy clara y conceptualmente correcta. El papel fisiológico del surfactante está perfectamente descrito, destacando su acción dinámica para igualar presiones. El cálculo de la diferencia de presión es correcto, con la fórmula y unidades adecuadas. & \textit{---} \\
    \midrule
    \textbf{Pregunta 3: Optimización en Hemodiálisis y Primera Ley de Fick} & 1.5/2.0 & La parte (a) sobre la determinación matemática de la tasa de transporte es correcta, aplicando bien la Ley de Fick y calculando el aumento de eficiencia. La parte (b) sobre el riesgo mecánico en la práctica clínica es también muy buena, identificando correctamente la ruptura de la membrana y sus consecuencias. Sin embargo, en la parte (c), aunque el cálculo inicial de la tasa de transporte en kg/s es correcto (1.8x10\textasciicircum{}-7 kg/s), la conversión posterior a otras unidades (indicada como 'x 100') es incorrecta, resultando en un valor erróneo (1.8x10\textasciicircum{}-5 kg/s). & \textit{Error en la conversión de unidades en el cálculo final de la tasa de transporte (parte c).} \\
    \midrule
    \textbf{Pregunta 4: Errores de Infusión Intravenosa y Ósmosis} & 1.0/2.0 & La parte (a) sobre el flujo neto de solvente y el fenómeno fisiopatológico con agua destilada es excelente, explicando detalladamente la hemólisis. Sin embargo, la parte (b) que solicita la reacción fisiológica al NaCl 5\% (solución hipertónica) está completamente ausente en su explicación cualitativa. En su lugar, el estudiante presenta el cálculo de la presión osmótica, que corresponde a la parte (c). El cálculo de la presión osmótica (parte c) es correcto. & \textit{Omisión de la explicación cualitativa de la reacción fisiológica al NaCl 5\% (parte b).; Confusión en la ubicación de la respuesta de la parte (c).} \\
    \midrule
    \textbf{Pregunta 5: Capilaridad y Toma de Muestras Sanguíneas} & 1.0/2.0 & Las partes (a) y (b) son correctas. El estudiante explica adecuadamente el comportamiento de la sangre en paredes hidrofóbicas (descenso capilar) y la ventaja de reducir el radio en microcapilares para el ascenso del líquido. Sin embargo, en la parte (c), el cálculo de la altura de ascenso es incorrecto (1.116 x 10\textasciicircum{}-5 metros) y no se muestra ningún procedimiento (fórmula, sustitución de valores), lo que impide identificar la fuente del error. & \textit{Cálculo numérico incorrecto de la altura de ascenso (parte c).; Ausencia de procedimiento y fórmula en el cálculo (parte c).} \\
    \midrule
\end{longtable}

% ══ ERRORES DETECTADOS ═══════════════════════════════════════════════════════
\section*{Análisis de errores}

\subsection*{Errores conceptuales}
\begin{itemize}[leftmargin=1.5em, itemsep=2pt]
  \item Aunque no es un error conceptual fundamental, la omisión de la explicación cualitativa sobre el efecto de una solución hipertónica (NaCl 5\%) en los eritrocitos (Pregunta 4b) sugiere una laguna en la comprensión completa de los fenómenos osmóticos en un contexto clínico.
\end{itemize}

\subsection*{Errores procedimentales}
\begin{itemize}[leftmargin=1.5em, itemsep=2pt]
  \item Error en la conversión de unidades al calcular la tasa de transporte en la Pregunta 3c.
  \item Falta de respuesta cualitativa para la Pregunta 4b.
  \item Cálculo numérico incorrecto y ausencia de pasos de procedimiento en la Pregunta 5c.
\end{itemize}

% ══ RECOMENDACIONES AL ESTUDIANTE ════════════════════════════════════════════
\section*{Retroalimentación para el estudiante}
\begin{tcolorbox}[cajaRecomendacion, title={Dirigido a: Daniela Vallenilla}]
Estimado estudiante, tu examen demuestra una sólida base en los principios de la física aplicada a tu campo. Has logrado explicar con claridad y precisión muchos conceptos importantes. Para mejorar aún más, te sugiero prestar especial atención a los detalles en los cálculos: siempre verifica tus conversiones de unidades y asegúrate de mostrar todos los pasos (fórmula, sustitución de valores y resultado final con unidades correctas). Además, es crucial que te asegures de responder todas las partes de cada pregunta, tanto las explicaciones cualitativas como los cálculos, para demostrar una comprensión completa del tema. Revisa los fenómenos osmóticos y sus efectos en las células con soluciones hipertónicas. ¡Sigue practicando y prestando atención a estos detalles para alcanzar un desempeño excelente!
\end{tcolorbox}

% ══ NOTA PARA EL PROFESOR ════════════════════════════════════════════════════
\section*{Nota para el profesor}
\begin{tcolorbox}[
  enhanced, breakable,
  colback=yellow!8, colframe=yellow!60!black,
  fonttitle=\bfseries, coltitle=black,
  title={Observaciones para revisión manual},
  top=4pt, bottom=4pt, left=8pt, right=8pt,
]
El estudiante presenta un buen nivel general. Los errores principales son de procedimiento (conversión de unidades, cálculo incorrecto sin procedimiento) y una omisión significativa en la Pregunta 4b (explicación cualitativa del NaCl 5\%). Las explicaciones conceptuales son, en su mayoría, muy buenas. El puntaje sugerido de 7.5/10 refleja un desempeño bueno, con margen de mejora en la rigurosidad de los cálculos y la completitud de las respuestas.
\end{tcolorbox}

% ══ PIE DE INFORME ════════════════════════════════════════════════════════════
\vspace{12pt}
\begin{center}
  \footnotesize\color{gray}
  Informe generado automáticamente por el sistema de evaluación con IA (gemini-2.5-flash) y soporte LaTex. \\
  Este documento es una evaluación \textbf{preliminar} para ser revisado por el profesor antes de ser entregado al 
  estudiante. \\
  Generado el 2026-08-08 \textbullet\ BIOANALISIS Sistema Académico de Evaluación.
  \end{center}

\end{document}
